% Source: https://artofproblemsolving.com/community/c3h1338626_help_please
% Stable identity: aops-topic-1338626-q01
Each angle adjacent to $x$ forms a straight line with it, so each adjacent angle is $180^\circ-x$.

\begin{asy}
usepackage("amsmath");
size(6cm);
defaultpen(fontsize(10pt));
pair O=(0,0);
pair a1=dir(25),a2=dir(155),a3=dir(205),a4=dir(-25);
draw(2*a3--2*a1,linewidth(0.9)); draw(2*a4--2*a2,linewidth(0.9));
draw(arc(O,0.45,25,155),blue); label("$x$",0.62*dir(90),blue);
draw(arc(O,0.7,155,205),red); label("$180^\circ-x$",0.95*dir(180),W,red);
draw(arc(O,0.7,-25,25),red); label("$180^\circ-x$",0.95*dir(0),E,red);
dot(O);
label("\mbox{}",(min(currentpicture,user=true)+max(currentpicture,user=true))/2,black); // keeps the text after the figure black when LaTeX typesets the labels inline
\end{asy}

The two adjacent angles are vertical angles, both equal to $180^\circ-x$, and their sum is $100^\circ$:
\[
2(180^\circ-x)=100^\circ\quad\Longrightarrow\quad 180^\circ-x=50^\circ\quad\Longrightarrow\quad x=\boxed{130^\circ}.
\]
(The answer $80^\circ$ would be the supplement of $100^\circ$, but $100^\circ$ is the sum of two angles, each of which is the supplement of $x$.)
